\section{Overview}

The experimental design evaluates both correctness and robustness. Correctness means that the methods behave consistently on clean datasets and match hand-computed fuzzy TOPSIS behavior in a small validation example. Robustness means that the methods can preserve the clean target rank when a subset of decision makers is contaminated.

\section{Dataset Types}

Three categories of data are used:

\begin{enumerate}
    \item Synthetic fuzzy group-decision datasets, generated with known decision-maker counts, criteria, alternatives, and contamination levels.
    \item Small fuzzy supplier-selection datasets, used as diagnostic examples.
    \item Real/pseudo-real datasets, where crisp datasets are converted into fuzzy group-decision panels for robustness testing.
\end{enumerate}

The main publication-facing datasets are:

\begin{itemize}
    \item Healthcare countries 2021: 26 alternatives and 13 criteria.
    \item Car evaluation: 300 alternatives and 6 criteria in the focused benchmark.
    \item Healthcare resource allocation: 300 alternatives and 18 criteria in the focused benchmark.
\end{itemize}

\section{Pseudo-Real Data Conversion}

Some datasets are not originally fuzzy group decision-making datasets. For these datasets, crisp values are normalized into a 1--9 rating scale and converted into triangular fuzzy numbers. Multiple pseudo decision makers are then created using bounded deterministic variation. This allows the robustness behavior of the methods to be tested on realistic data structures while preserving a controlled experimental setting.

This conversion is not presented as original human expert data. It is a pseudo-real simulation protocol for testing algorithmic robustness. The thesis therefore clearly distinguishes between native fuzzy group data and pseudo-real converted panels.

\section{Contamination Protocol}

The contamination protocol selects the clean lowest-ranked alternative as the attack target. A fraction of decision makers is then contaminated. Contaminated decision makers assign high fuzzy ratings to the target and low fuzzy ratings to the other alternatives. This produces a structured target-promotion attack.

Attack fractions are tested at 10\%, 20\%, 30\%, 40\%, 50\%, and 60\%. Because the number of decision makers is discrete, the effective attack fraction may differ slightly from the requested fraction. For example, with 15 decision makers, four contaminated decision makers represent 26.7\% effective contamination.

\section{Evaluation Metrics}

The following metrics are used:

\begin{itemize}
    \item Spearman rank correlation against the clean M1 reference.
    \item Kendall rank correlation against the clean M1 reference.
    \item Top-1 agreement with the clean reference.
    \item Target rank under contamination.
    \item Target-rank error.
\end{itemize}

Target-rank error is the most important metric for this thesis. In targeted manipulation, global rank correlation can remain high even when the attacked target moves from the bottom of the ranking to the top. Therefore, a method can look good under Spearman or Kendall correlation while still failing the actual security objective.

\section{Repeated Runs}

Stochastic methods use deterministic seeds and repeated runs. The focused benchmark uses 30 repeats and 200 bags. Deterministic methods are run once where appropriate. Repeated runs allow standard deviations, confidence intervals, and paired target-error comparisons to be reported.

\section{External Comparator Baselines}

Five external comparator baselines are included:

\begin{itemize}
    \item EB1 Median TOPSIS: component-wise median aggregation.
    \item EB2 Trimmed-Mean TOPSIS: component-wise 20\% trimmed mean aggregation.
    \item EB3 MAD-Consensus TOPSIS: filters decision makers far from the robust median vector.
    \item EB4 Individual-Borda TOPSIS: runs fuzzy TOPSIS per decision maker and aggregates rankings by Borda voting.
    \item EB5 Huang--Li Group-Ideal TOPSIS: adapts the Huang--Li group-ideal TOPSIS aggregation model using individual fuzzy TOPSIS closeness scores.
\end{itemize}

\section{Implementation Environment}

All methods are implemented in Python. The core scripts include:

\begin{itemize}
    \item \texttt{m1\_normal\_topsis.py}
    \item \texttt{m2\_bagged\_topsis.py}
    \item \texttt{m3\_ml\_filtered.py}
    \item \texttt{m4\_weighted\_bagging.py}
    \item \texttt{m5\_cluster\_stratified.py}
    \item \texttt{m6\_reliability\_weighted.py}
    \item \texttt{m7\_entropy\_reliability.py}
\end{itemize}

The final evidence tables are generated from saved benchmark outputs using \texttt{compile\_final\_evidence.py}. Statistical summaries are generated using \texttt{analyze\_statistics.py}.
